\documentclass[10pt]{beamer}
\usepackage[utf8]{inputenc}
\usepackage[T1,T2A]{fontenc}
\usepackage[russian]{babel}
\usepackage{color}
\usepackage{calc}
\usepackage{graphicx}
\usepackage{epstopdf}
\usepackage{hyperref}
\hypersetup{unicode,colorlinks}
\usepackage{csquotes}
\usepackage{upquote}
\usepackage{cprotect}
\usetheme[progressbar=head,numbering=fraction,block=fill]{metropolis}
\usepackage{dejavu}
\usepackage{etoolbox}
\usepackage{bm}
\usepackage[ND]{prftree}
\usepackage[tableaux]{prooftrees}
\usepackage{mathtools} % for bigtimes
\usepackage{diagbox}
\usepackage{verbatimbox}
\usepackage{fitch}
% fitch по умолчанию печатает эти правила как $\Rightarrow$I/E; в курсе импликация всюду $\to$.
\gappto\ndrules{\def\ii{\by{$\to$I}}\def\ie{\by{$\to$E}}}
%\usepackage[defaultsans]{droidsans}
%\usepackage[defaultmono]{droidmono}
%\makeatletter
%\input droid/t1fdm.fd
%\makeatother

\makeatletter
\DeclareRobustCommand*{\bfseries}{%
    \not@math@alphabet\bfseries\mathbf
    \fontseries\bfdefault\selectfont
    \boldmath
}
\makeatother

\makeatletter
\chardef\straightquote@code=\catcode`'
\chardef\backquote@code=\catcode``
\catcode`'=\active \catcode``=\active
\patchcmd{\@noligs}
{\textasciigrave}
{\fixedtextasciigrave}
{}{}
\newcommand{\fixedtextasciigrave}{%
    \makebox[.5em]{\fontencoding{TS1}\fontfamily{fvs}\selectfont\textasciigrave}% Vera Sans
}
\catcode\lq\'=\straightquote@code
\catcode\lq\`=\backquote@code
\makeatother

\setbeamercovered{again covered={\opaqueness<1->{25}}}

\forestset{close with=\ensuremath{\bigtimes}}

\newcommand{\customframefont}[1]{
    \setbeamertemplate{itemize/enumerate body begin}{#1}
    \setbeamertemplate{itemize/enumerate subbody begin}{#1}
}

\NewEnviron{framefont}[1]{
    \customframefont{#1} % for itemize/enumerate
    {#1 % For the text outside itemize/enumerate
        \BODY
    }
    \customframefont{\normalsize}
}


\makeatletter
\DeclareRobustCommand{\rvdots}{%
    \vphantom{\int\limits^x}\smash[t]{\vdots}
}
\makeatother

% Короткое вертикальное многоточие: размеры в em, поэтому оно уменьшается вместе со шрифтом
% (у \vdots они фиксированы в pt).
\DeclareRobustCommand{\ndvdots}{%
    \mathord{\vbox{\baselineskip=.3em\lineskiplimit=0pt\kern.15em\hbox{.}\hbox{.}\hbox{.}}}}
% Подвывод в схеме правила, нарисованный как во fitch: вертикальная черта и флажок под допущением.
% Выравнивание по нижней строке, чтобы в \lor E заключения подвыводов стояли на одной линии с A \lor B.
\newcommand{\ndsubproof}[2]{{%
    \setlength{\arrayrulewidth}{.2mm}%
    \begin{array}[b]{@{}|@{\,}l@{\,}}#1\\\hline\ndvdots\\#2\end{array}}}

\newcommand{\N}{\mathbb{N}}
\newcommand{\Z}{\mathbb{Z}}
\newcommand{\Q}{\mathbb{Q}}
\newcommand{\R}{\mathbb{R}}
\newcommand{\pow}[1]{\mathcal{P}({#1})}
\newcommand{\continuum}{\mathfrak{c}}

% \reveal<спец>{...}: показать содержимое с указанного слайда, до этого оставив пустое место той же ширины.
% В отличие от \pause и \uncover не оставляет невидимых PGF-специалов; внутри выключной формулы \pause
% оставляет их в пустом абзаце после формулы, и TeX набирает лишнюю пустую строку.
% Следующий \item после такой формулы получает <.(1)-> (тот же слайд, что и \reveal<+(1)->)
% или <+(1)-> (следующий слайд); в beamer «.» "--- это значение счётчика минус один.
\newcommand<>{\reveal}[1]{\alt#2{#1}{\phantom{#1}}}

% Многобуквенные имена в формулах: курсив с нормальным кернингом (а не произведение букв).
\newcommand{\ident}[1]{\mathit{#1}}
\newcommand{\Var}{\ident{Var}}
\newcommand{\Prop}{\ident{Prop}}
\newcommand{\Const}{\ident{Const}}
\newcommand{\Fun}{\ident{Fun}}
\newcommand{\Pred}{\ident{Pred}}
\newcommand{\Term}{\ident{Term}}
\newcommand{\Frm}{\ident{Form}} % \Form занят hyperref
\newcommand{\arity}{\ident{arity}}
\newcommand{\Th}{\ident{Th}}
\newcommand{\Dom}{\ident{Dom}}
\newcommand{\Ran}{\ident{Ran}}
\newcommand{\id}{\ident{id}}
\newcommand{\Ord}{\ident{Ord}}
\newcommand{\RAA}{\ident{RAA}}
% Названия теорий и аксиом: прямой шрифт.
\newcommand{\ZF}{\mathrm{ZF}}
\newcommand{\ZFC}{\mathrm{ZFC}}
\newcommand{\AC}{\mathrm{AC}}
\newcommand{\DC}{\mathrm{DC}}
\newcommand{\CH}{\mathrm{CH}}
\newcommand{\GCH}{\mathrm{GCH}}

\author{Алексей Романов}
\subtitle{Математическая логика и теория алгоритмов}
%\logo{}
\institute{МИЭТ}
\subject{Математическая логика и теория алгоритмов}
%\setbeamercovered{transparent}
%\setbeamertemplate{navigation symbols}{}


% Скрываем обоснование целиком, сохраняя его место в доказательстве.
\newcommand{\byon}[3]{\by{\uncover<#1>{#2, \ndref{#3}}}{}}

\title{Натуральная дедукция (естественный вывод)}
\date{\today}

\begin{document}
\begin{frame}[plain]
    \maketitle
\end{frame}

\begin{frame}
    \frametitle{Секвенции}
    \begin{itemize}
        \item Выражение $A_1,\ldots,A_n \vdash B$ называется секвенцией и читается \enquote{из $A_1,\ldots,A_n$ выводится $B$}.
        \item Слева --- множество посылок: порядок и повторения неважны. Оно может быть пустым. \pause
        \item Секвенция \emph{общезначима}, если на всех наборах, где посылки истинны, $B$ тоже истинна. \pause
        \item При $n\geqslant1$ это равносильно тождественной истинности $A_1\land\cdots\land A_n\to B$.
        \item Без посылок: $\vdash B$ общезначима, если $B$ тождественно истинна. \pause
        \item $A\equiv B$ тогда и только тогда, когда секвенции $A\vdash B$ и $B\vdash A$ общезначимы.
    \end{itemize}
\end{frame}

\begin{frame}
    \frametitle{Правила натуральной дедукции}
    \tiny
    \begin{itemize}
        \item Для $\land$:
        \[ 
        \frac{A \qquad B}{A \wedge B}\ \wedge I 
        \qquad \qquad 
        \frac{A \wedge B}{A}\ \wedge{E}
        \qquad
        \frac{A \wedge B}{B}\ \wedge{E} 
        \]
        \item Для $\to$:
        \[ 
        \frac{\boxed{\begin{matrix}
                A \\
                \rvdots \\
                B
        \end{matrix}}}{A \to B}\ \to I 
        \qquad \qquad 
        \frac{A \to B \quad A}{B}\ \to E
        \]
        \item Для $\neg$ и $\bot$:
        \[
        \frac{\boxed{\begin{matrix}
                A \\
                \rvdots \\
                \bot
        \end{matrix}}}{\neg A}\ \neg I 
        \qquad \qquad 
        \frac{\neg A \quad A}{\bot}\ \neg E/\bot I
        \qquad \qquad 
        \frac{\bot}{A}\ \bot{E}
        \]
        \item Для $\lor$:
        \[
        \frac{A}{A \lor B}\ \lor{I}
        \qquad
        \frac{B}{A \lor B}\ \lor{I}
        \qquad \qquad 
        \frac{
            \begin{matrix}
                \phantom{A} \\
                \phantom{\rvdots} \\
                A \lor B 
            \end{matrix}
            \quad
            \boxed{\begin{matrix}
                A \\
                \rvdots \\
                C
            \end{matrix}}
            \quad
            \boxed{\begin{matrix}
                B \\
                \rvdots \\
                C
            \end{matrix}}
        }{C}\ \lor E
        \qquad \qquad         
        \frac{\boxed{\begin{matrix}
                \neg A \\
                \rvdots \\
                B
        \end{matrix}}}{A \lor B}\ \lor I'
        \qquad 
        \frac{\boxed{\begin{matrix}
                \neg B \\
                \rvdots \\
                A
        \end{matrix}}}{A \lor B}\ \lor I'
        \]
        \item Остальные:
        \[
        \frac{\boxed{\begin{matrix}
                \neg A \\
                \rvdots \\
                \bot
        \end{matrix}}}{A}\ \RAA
        \qquad \qquad 
        \frac{A}{A}\ R
        \]
        \item $I$ обозначает правила введения (Introduction), $E$ "--- правила исключения (Elimination). 
        \item[] Например, $\land I$ это правило \enquote{введения конъюнкции}. Можете также писать как В$\land$.
    \end{itemize}
\end{frame}

\begin{frame}
    \frametitle{Выводы и подвыводы}
    \begin{itemize}
        \item Строка вывода: формула и обоснование (правило, номера строк или начала и конца подвывода).
        \item Подвывод начинается с допущения и отмечается дополнительной вертикальной чертой. \pause
        \item Можно ссылаться на более ранние строки, кроме строк уже закрытых подвыводов.
        \item Правило может использовать закрытый подвывод целиком, снимая его допущение. Остальные посылки и допущения остаются. \pause
        \item Секвенция \emph{выводима}, если есть вывод $B$ из посылок по правилам. \pause
        \item \textbf{Теорема.} В классической натуральной дедукции секвенция выводима тогда и только тогда, когда общезначима.
    \end{itemize}
\end{frame}

\begin{frame}
    \frametitle{Поиск доказательства}
    \begin{itemize}
        \item Есть одна или несколько целей со своими посылками, допущениями и уже полученными формулами. \pause
        \item Выбираем цель и правило, по которому можно её получить (часто правило введения главной связки).
        \item Подвывод открываем, если выбранное правило этого требует.
        \item Ещё не полученные посылки правила (или цели подвывода) становятся новыми целями. \pause
        \item Например:
        \begin{itemize}
            \item Для $(p\land q)\land r$ по $\land I$ цели: $p\land q$ и $r$.
            \item Для $\neg p\to\neg q$ по $\to I$: подвывод с допущением $\neg p$ и целью $\neg q$.
        \end{itemize} \pause
        \item Если цель уже доступна, используем её по правилу $R$.
        \item Если путь не удаётся, пробуем другое правило.
        \item Доказательство готово, когда достигнуты все цели.
    \end{itemize}
\end{frame}

\begin{frame}
    \frametitle{Пример}
    \begin{align*}
        \begin{nd}
            \hypo{1}{(p \land q) \land r} \by{Дано}{} \onslide<1->
            \have{2}{p \land q} \ae{1} \onslide<4->
            \have[n][-4]{5}{q} \byon{4-}{$\land E$}{2} \onslide<3->
            \have[n][-3]{3}{r} \ae{1} \onslide<3->
            \have[n][-2]{4}{p} \byon{4-}{$\land E$}{2} \onslide<2->
            \have[n][-1]{6}{q \land r} \byon{3-}{$\land I$}{5,3} \onslide<2->
            \have{n}{p \land (q \land r)} \byon{2-}{$\land I$}{4,6} \onslide<1->
        \end{nd}
    \end{align*}
\end{frame}

\begin{frame}
    \frametitle{Пример}
    \begin{align*}
        \begin{nd}
            \open
            \onslide<2->    \hypo{prem}{p} \by{Дано}{}
                \open
            \onslide<3->        \hypo{hyp}{q} \by{Дано}{}
            \onslide<3->    \have{reit}{p} \r{prem}
                \close
            \onslide<2->    \have[n][-1]{cond}{q \rightarrow p} \byon{3-}{$\to I$}{hyp-reit}
            \close
            \onslide<1-> \have[n]{conc}{p \rightarrow (q \rightarrow p)} \byon{3-}{$\to I$}{prem-(cond)}
        \end{nd}
    \end{align*}
\end{frame}

\begin{frame}
    \frametitle{Пример: транзитивность импликации}
    \begin{itemize}
        \item Докажем $p\to q,\ q\to r\vdash p\to r$.
        \item<2-> Для введения $\to$ предположим $p$ и выведем $r$.
    \end{itemize}
    {\footnotesize
    \begin{align*}
        \begin{nd}
            \hypo{t1}{p\to q} \by{Дано}{} \onslide<1->
            \hypo{t2}{q\to r} \by{Дано}{} \onslide<1->
            \open
                \hypo{t3}{p} \by{Допущение}{} \onslide<2->
                \have{t4}{q} \ie{t1,t3} \onslide<4->
                \have{t5}{r} \byon{4-}{$\to E$}{t2,t4} \onslide<3->
            \close
            \have{t6}{p\to r} \byon{4-}{$\to I$}{t3-t5} \onslide<1->
        \end{nd}
    \end{align*}
    }
    \onslide<5->
    \begin{itemize}
        \item Допущение $p$ снято; посылки $p\to q$ и $q\to r$ остались.
    \end{itemize}
\end{frame}

\begin{frame}
    \frametitle{Пример: введение отрицания}
    \begin{itemize}
        \item Докажем $p\to q,\ \neg q\vdash\neg p$.
        \item<2-> Чтобы получить $\neg p$, предположим $p$ и выведем $\bot$.
    \end{itemize}
    {\footnotesize
    \begin{align*}
        \begin{nd}
            \hypo{n1}{p\to q} \by{Дано}{} \onslide<1->
            \hypo{n2}{\neg q} \by{Дано}{} \onslide<1->
            \open
                \hypo{n3}{p} \by{Допущение}{} \onslide<2->
                \have{n4}{q} \ie{n1,n3} \onslide<4->
                \have{n5}{\bot} \byon{4-}{$\neg E$}{n2,n4} \onslide<3->
            \close
            \have{n6}{\neg p} \byon{4-}{$\neg I$}{n3-n5} \onslide<1->
        \end{nd}
    \end{align*}
    }
    \onslide<5->
    \begin{itemize}
        \item Снимая допущение $p$, получаем $\neg p$.
    \end{itemize}
\end{frame}

\begin{frame}
    \frametitle{Пример: снятие двойного отрицания}
    \begin{itemize}
        \item Докажем $\neg\neg p\vdash p$ с помощью $\RAA$.
        \item<2-> Предположим отрицание цели: $\neg p$.
    \end{itemize}
    {\footnotesize
    \begin{align*}
        \begin{nd}
            \hypo{r1}{\neg\neg p} \by{Дано}{} \onslide<1->
            \open
                \hypo{r2}{\neg p} \by{Допущение}{} \onslide<2->
                \have{r3}{\bot} \by{$\neg E$}{r1,r2} \onslide<3->
            \close
            \have{r4}{p} \byon{3-}{$\RAA$}{r2-r3} \onslide<1->
        \end{nd}
    \end{align*}
    }
    \onslide<4->
    \begin{itemize}
        \item $\neg I$ для $2\!\!-\!\!3$ дало бы $\neg\neg p$. Именно $\RAA$ позволяет заключить $p$.
    \end{itemize}
\end{frame}

\begin{frame}
    \frametitle{Пример: разбор случаев}
    \begin{itemize}
        \item Докажем $p\lor q,\ \neg p\vdash q$.
    \end{itemize}
    {\footnotesize
    \begin{align*}
        \begin{nd}
            \hypo{d1}{p\lor q} \by{Дано}{} \onslide<1->
            \hypo{d2}{\neg p} \by{Дано}{} \onslide<1->
            \open
                \hypo{d3}{p} \by{Допущение}{} \onslide<2->
                \have{d4}{\bot} \by{$\neg E$}{d2,d3} \onslide<4->
                \have{d5}{q} \byon{4-}{$\bot E$}{d4} \onslide<3->
            \close
            \open
                \hypo{d6}{q} \by{Допущение}{} \onslide<5->
                \have{d7}{q} \r{d6} \onslide<5->
            \close
            \have{d8}{q} \byon{5-}{$\lor E$}{d1,d3-d5,d6-d7} \onslide<1->
        \end{nd}
    \end{align*}
    }
    \onslide<6->
    \begin{itemize}
        \item $q$ получено в обоих случаях; оба допущения сняты.
    \end{itemize}
\end{frame}

\begin{frame}
    \frametitle{Пример: коммутативность дизъюнкции}
    \begin{itemize}
        \item Докажем $p\lor q\vdash q\lor p$.
        \item<2-> Разберём случаи $p$ и $q$. В каждом применим $\lor I$.
    \end{itemize}
    {\footnotesize
    \begin{align*}
        \begin{nd}
            \hypo{c1}{p\lor q} \by{Дано}{} \onslide<1->
            \open
                \hypo{c2}{p} \by{Допущение}{} \onslide<2->
                \have{c3}{q\lor p} \byon{3-}{$\lor I$}{c2} \onslide<3->
            \close
            \open
                \hypo{c4}{q} \by{Допущение}{} \onslide<4->
                \have{c5}{q\lor p} \byon{5-}{$\lor I$}{c4} \onslide<5->
            \close
            \have{c6}{q\lor p} \byon{5-}{$\lor E$}{c1,c2-c3,c4-c5} \onslide<1->
        \end{nd}
    \end{align*}
    }
    \onslide<6->
    \begin{itemize}
        \item В обоих случаях получена одна цель $q\lor p$. Оба допущения сняты.
    \end{itemize}
\end{frame}

\begin{frame}
    \frametitle{Допустимые правила}
    \begin{itemize}
        \item Выше доказали $(p \land q) \land r \vdash p \land (q \land r)$.
        \item Это можно превратить в доказательство $(A \land B) \land C \vdash A \land (B \land C)$ для любых формул $A,B,C$. Как? \pause
        \item Одновременно заменим $p,q,r$ на $A,B,C$ во всём доказательстве. \pause
        \item Это даёт новое \emph{допустимое правило}:
        \[
            \frac{(A \land B) \land C}{A \land (B \land C)}\ \wedge A
        \]
        \item И так для каждой секвенции, которую докажем! Например
        \[
            \frac{\neg \neg A}{A}\ \neg \neg E \qquad \qquad \frac{A\lor B \qquad \neg A}{B}\ \ident{DS} \qquad \qquad \frac{A \to B \qquad B \to C}{A \to C}\ \to T
        \]
        и т.д.
        \item У каких-то из этих правил есть обозначения, как выше, но можно просто пометить соответствующей секвенцией.
    \end{itemize}
\end{frame}

\begin{frame}
    \frametitle{Допустимость $\lor I'$}
    \begin{itemize}
        \item Дан подвывод $\Pi$: из допущения $\neg B$ получена $A$.
        \item<5-> В строке 6 вставляем $\Pi$, используя $\neg B$ из строки 5.
    \end{itemize}
    {\footnotesize
    \begin{align*}
        \begin{nd}
            \open
                \hypo{v1}{\neg(A\lor B)} \by{Допущение}{} \onslide<2->
                \open
                    \hypo{v2}{B} \by{Допущение}{} \onslide<4->
                    \have{v3}{A\lor B} \byon{4-}{$\lor I$}{v2} \onslide<4->
                    \have{v4}{\bot} \byon{4-}{$\neg E$}{v1,v3} \onslide<4->
                \close
                \have{v5}{\neg B} \byon{4-}{$\neg I$}{v2-v4} \onslide<3->
                \have{v6}{A} \byon{5-}{$\Pi$}{v5} \onslide<5->
                \have{v7}{A\lor B} \byon{6-}{$\lor I$}{v6} \onslide<6->
                \have{v8}{\bot} \byon{6-}{$\neg E$}{v1,v7} \onslide<6->
            \close
            \have{v9}{A\lor B} \byon{6-}{$\RAA$}{v1-v8} \onslide<1->
        \end{nd}
    \end{align*}
    }
\end{frame}

\end{document}
